\section{A closed-gap source criterion for FULL A2 under a uniform window}
Write $[t]=1+\cdots+z^{t-1}$. All coefficient sequences are zero extended.
The general original conjecture remains OPEN.

\begin{theorem}[A complete source-to-window A2 bridge]
Let $p$ be an actual old ordinary arm product. Assume FULL OLD A2
and coefficient tilt A1 for $p$. Put $u=(\Delta p)^+$, let $m$ be
the LAST ordinary mode of $p$, and let $\ell$ be the first mode of
$2^nu_n$. Set $E_n=2p_n-p_{n-1}$.
Assume the following CLOSED-gap proper comparison for every $t\ge1$:
whenever $n>m$, $k=n-t\ge0$, $k<\ell$,
\[
 p_n>p_k,\quad p_{n+1}>p_{k+1},\quad 0\le E_n\le E_k,
\]
one has $E_{n+1}-E_n\le E_{k+1}-E_k$.
Then EVERY ordinary uniform window $[t]p$ has FULL A2.
This theorem concludes weighted-positive-D UI, not whole raw-source UI.
\end{theorem}
\begin{proof}
Put $v=[t]p$ and $w_n=\Delta v_n=p_n-p_{n-t}$. Uniform convolution
makes $v$ ordinarily UI. Its positive differences have initial interval
support: before the old first mode the entering term strictly exceeds
the outgoing term; across the old mode interval its differences are
nonincreasing; after that interval they are nonpositive. This also
follows directly from the established actual positive-difference support.
Write $\sigma$ for its first mode and $f_n=2^n(w_n)^+$.

First build a COMPLETE UI head comparison, including its first boundary.
For $n\le m$, every signed old difference is nonnegative, so
$w_n=([t]u)_n$. Its weighted sequence is convolution of the old A2
sequence by $(1,2,\ldots,2^{t-1})$, a log-concave kernel without holes.
Such convolution preserves UI. Denote this WHOLE UI sequence by $g$.
At $m+1$, signed convolution and clipping give $f_{m+1}\le g_{m+1}$.
Alternatively, for $n<t$ one has $w_n=p_n$ and $f_n=2^np_n$,
whose whole sequence is UI by A1; at $t$, $f_t\le2^tp_t$.
Choose
\[
                         c=\max(m,t-1).
\]
If $m\ge t-1$ use the first head sequence; otherwise use the second.
In either case $f=g$ through $c$, $f_{c+1}\le g_{c+1}$, and $g$ is
WHOLE UI. An earlier strict fall of $f$ therefore forces
$f_{c+1}\le f_c$. It remains to propagate NONSTRICT falls after this
boundary; using strict falls alone would leave an equality gap.

At a crossing $n>m$, $n\ge t$, if $w_n,w_{n+1}>0$, set $k=n-t$.
If $k\ge\ell$, the old weighted geometric tail gives
\[
 w_n\le\sum_{i=k+1}^{\mathrm{firstmode}(p)}u_i\le2u_{k+1}.
\]
The finite sum bound follows from $2u_{i+1}\le u_i$ throughout the
used tail; its sharper telescoping form subtracts the last coefficient.
Also
\[
 w_{n+1}=w_n-u_{k+1}-(p_n-p_{n+1}),\qquad p_n-p_{n+1}\ge0,
\]
so $2w_{n+1}\le w_n$ unconditionally in this branch.

If $k<\ell$, then $E_k,E_{k+1}>0$ and
$H=E_{k+1}-E_k=2u_{k+1}-u_k\ge0$. A nonstrict weighted fall means
\[
          2w_n-w_{n-1}=E_n-E_k\le0.                    \tag{1}
\]
If $E_n<0$, A1 implies $E_{n+1}\le0$, hence
$2w_{n+1}-w_n=E_{n+1}-E_{k+1}<0$.
Otherwise the CLOSED-gap hypothesis and (1) give
\[
 E_{n+1}-E_{k+1}\le E_n-E_k\le0,
\]
again proving $2w_{n+1}\le w_n$. Once the positive-difference support
ends, $f$ is zero and every later edge is nonincreasing.
Thus nonstrict falls propagate throughout the entire crossing range.
The kernel-checked head/boundary/closed-cross lemma excludes every
strict fall followed by a later rise of $f$. Finiteness supplies a
whole up/down mode. This proves FULL A2 for $[t]p$.
\end{proof}

\paragraph{The relative raw-kernel theorem retains a strict surplus at equality.}
The frozen relative-head proof assumes $0\le E_n<E_k$, but its exact
lower bound is stronger: with $\lambda=C/d_r>0$,
\[
 H\ge\lambda E_k+1,\qquad E_n\ge d_r>0,\quad T\le C.
\]
Thus $E_n\le E_k$ already gives $H\ge C+1>T$. Its proof of the
raw-kernel convolution inequality used neither strictness of the gap
nor C. The CLOSED-gap extension is therefore valid, with no appeal to
an unproved continuity or equality argument. It covers every $r\ge k$.

\section{First-missing-nine FULL A2 and ORIGINAL unamplified coefficients}
\begin{theorem}
Every final-deletion-closed class of ACTUAL normalized increasing arms
with FULL C and $\kappa\le9$ has FULL pair A2. The ORIGINAL strict
core and BOTH relatives are coefficient UI at EVERY strict peak,
including ONE arbitrary legal heavier append.
For an old nonrectangular pair with maximum first missing weight nine,
EVERY ordinary uniform window of its reverse product has FULL A2.
No whole raw-source-window UI conclusion is asserted.
\end{theorem}

\subsection{All positive mixed recoveries have shared forward rank at most ten}
The old maximum-$\kappa\le8$ class and arbitrary heavier append are
frozen. Consider an old pair of maximum first missing weight nine,
with FULL OLD A2 in the count induction. A1 follows from FULL C.
Both-rising C and both-nonrising actual resources give $T\le0$ after
the LAST ordinary product mode whenever the current source is nonnegative.
The latter includes both halves by combining the unconditional forward
both-nonincreasing theorem and the actual reverse half-tail theorem.
Only mixed recovery remains.

If its rising opposite has $\kappa\le8$, the complete frozen pointwise
exceptions have owner weight at most eight. An owner of first missing
weight nine has weight at least 36, so none is possible. Hence a positive
new mixed recovery must have opposite first missing weight NINE.
With $Z=1^\alpha U$ the nonincreasing owner, use the exact initial
exchange $\alpha_* =\max(\alpha,t_1-1)$.
The frozen absolute-nine theorem pays all owners with $h\ge6$.
The gap-six theorem pays $h=4,5$ when $\alpha_*\ge3$; the gap-five
theorem pays $h=2,3$ when $\alpha_*\ge4$.
The remaining canonical domain is
\[
\begin{array}{c|c|c|c}
 \alpha& h& t_1& D_{\max}\\\hline
 1&2,3,4,5&2&71\\
 2&2,3,4,5&2,3&56\\
 3&2,3&2,3,4&36
\end{array}                                                     \tag{2}
\]
Here $D=\sum(t_i-1)$, $d=9-\alpha$, and the unbounded necessary
bound is $D\le\min(2d(d-3),12d-25)$.

For $r=a-j\ge0$, positive recovery forces positive primitive curvature,
$D=2r+2+e$, and $0\le e\le d-3$. Put $L=d-e-3$ and
\[
 R_m=\frac{\binom{r+e+m+2}{m}}{\binom r m}.
\]
The genuine compulsory source increments give the frozen sufficient
rank payment
\[
 \frac{\mathbf1_{L\le r}}{R_L}
 +\sum_{\substack{2\le v\le6\\L+v\le r}}
       \frac{2^{v-2}}{R_{L+v}}\ \ge\ \frac{e+1}{r+1}.    \tag{3}
\]
All these offsets are in the mode-free early-source range through
$\kappa-3=6$; included birth reads have nonnegative actual indices.
The old CI-relative/reflection proof of (3) applies for every $h\ge2$,
not only the $h\ge6$ specialization of the earlier nine-threshold result.

The COMPLETE rank grammar (2)--(3) has 406 states. Exact rational
arithmetic pays 291. Its remaining 115 states have $D\le25$.
For EVERY remaining state enumerate all sorted partitions of $D$ into
the specified counts, with minimum increment at most $\alpha$.
This reconstructs 2143 distinct canonical owners and ALL 4377 used
owner/rank inputs. Of these, 3757 have nonpositive curvature and are
irrelevant to positive recovery. At the 620 positive-curvature inputs,
the direct complete compulsory-source bound pays 601:
\[
 F_Z(v)+F_Z(v-2)+2F_Z(v-3)+4F_Z(v-4)
                  +8F_Z(v-5)+16F_Z(v-6)\ge C,
 \qquad v=a-r-8.                                      \tag{4}
\]
The ONLY 19 source-unpaid candidates all have $j=a-r\le10$.
All rational numerators/positive denominators, nonpositive-curvature
classifications, payments (4), and candidate rank bounds reduce to
719 distinct nonnegative integer obligations in six ordinary-kernel
chunks and ONE complete public aggregate. An independent degree/maximum-part
grammar must reconstruct the entire rank and owner partition, every
category, every numerator, and the public obligations.

For $r\le-2$ no primitive curvature is available. At $r=-1$, the
unit birth is absent only when $a+2-9<0$, so $a\le6$ and $j=a+1\le7$.
One-part owners have positive primitive curvature ONLY at this same
negative index and are handled by the identical bound. Therefore
\[
             T>0\quad\Longrightarrow\quad j\le10       \tag{5}
\]
throughout the entire new mixed/source-prefix region, independently
of every opposite capacity, count, padding and weight.

\subsection{A genuine strict opposite head pays all nineteen candidates}
At a proper crossing, BOTH strict positive windows give
\[
                  k+n+2<a+b,\qquad k+1\le j-2\le8.
\]
This uses actual upper-centre reflection of $q_Zq_Y$, not a synthetic
scalar rank profile. Also $j\le a+1$, hence $k+1\le a-1$.
The entire used RAW owner source $R_Z=(2-z)q_Z$ is consequently
nonnegative, because it agrees with the corrected positive source
below degree $a$; its constant coefficient is two.

The opposite main interval product has degree $M$ and largest factor
degree $B$. Its complete first-missing staircase forces
\[
 M\ge1+\frac{7\cdot8}{2}=29,\qquad
 M-B\ge1+\frac{6\cdot7}{2}=22.
\]
The exact CI positive-primitive support formula gives positive
differences through $\min(\lfloor M/2\rfloor,M-B)\ge14$.
Every other actual shifted-CI summand has endpoint sum at least $M$,
so its differences are nonnegative through $\lfloor M/2\rfloor$.
Thus the ACTUAL opposite coefficients satisfy
\[
                       \Delta q_Y(i)\ge1\quad(0\le i\le14).      \tag{6}
\]
No coefficientwise lower bound is differentiated to obtain (6).

The exact head increment is $H=[z^{k+1}]R_Z(1-z)q_Y$.
All used terms are nonnegative. If $0\le r<k$, its owner term $r+1$
and (6) give
\[
 H\ge R_Z(r+1)\Delta q_Y(k-r)
   \ge R_Z(r+1)=q_Z(r+1)+d_{r+1}\ge C\ge T.
\]
For $r=-1$ the owner-constant term gives $H\ge2\ge T$.
For $r\ge k$, the CLOSED-gap relative raw-kernel proof above gives
$H>T$. Every positive recovery is therefore paid at EVERY proper
closed-gap head. For $T\le0$, FULL OLD A2 gives $H\ge0$.
Terminal forward ranks zero/one are paid directly or have negative
current source. This proves the complete closed-gap source criterion.

\subsection{FULL induction and original coefficients}
The first theorem gives FULL A2 for EVERY uniform window of an old
maximum-nine pair. In the non-one-count induction orient a heavier
new arm $Zt$. Empty, rectangular and heavier-one-part bases are frozen.
If the old pair has maximum first missing weight at most eight, its
frozen arbitrary heavier-append theorem applies. Otherwise use the new
window theorem. The exact actual correction starts at the new heavier
weight, and the NEW first ordinary mode is strictly below that weight.
Hence its positive differences are a terminal truncation of the window's
positive differences, followed by zeros; FULL A2 is preserved.
This closes the induction and also proves ONE arbitrary extra legal
heavier append without requiring C of that new arm. The frozen proper
deletion and lower-peak/reflection joins then give ORIGINAL coefficient
UI of the strict core and BOTH relatives.

Let A15 be A14 together with ALL max-four-prefix arms followed by FOUR
successive dominant capacities $T_i\ge2(W_0+\sum_{j<i}T_j)$, with
$T_1\ge\max(5,2W_0)$, and all max-eight tails of nonbinary weight at
most 32. Binary counts and initial padding are arbitrary. Frozen
dominant C and max-four/fixed-weight C give FULL C; final deletion
closure and $\kappa\le9$ are exact. Thus BOTH original arms may
independently have FOUR such jumps, arbitrary strict peak, and ONE
further legal heavier append. No amplification is used.

\paragraph{Verification scope and present status.}
NINE new abstract standard-axiom reports check CLOSED source propagation,
the negative tail, whole finite geometric budgets, the relative surplus
at equality, the WHOLE UI-head/boundary/closed-cross no-revival proof,
head selection and complete proper-head guards. ONE complete finite
aggregate checks all 719 distinct obligations. The actual convolution,
CI, old A1/A2/source and original-model assemblies above are written
proofs, NOT an end-to-end Lean composition implementation. The complete
406-rank/2143-owner finite premise requires independent reconstruction.
Additional actual windows and greedy originals are QA ONLY. General C,
general unbounded-first-missing A2, and the general ORIGINAL conjecture
remain OPEN. No whole raw-source-window shape is inferred from the
new weighted-positive-difference theorem.
